% ============================================================
%  packages.tex  --  alle eingebundenen LaTeX-Pakete
%
%  WARUM KEIN \usepackage{babel} MIT NGERMAN?
%  ------------------------------------------
%  Babel ist das Standard-Lokalisierungspaket fuer LaTeX und
%  liefert u.a. deutsche Trennmuster, Anpassung von
%  Kapitelbeschriftungen ("Kapitel" statt "Chapter") und
%  korrekte Silbentrennung.
%
%  Auf diesem System (TeX Live 2023/Debian) fehlt das Paket
%  "babel-german" (normalerweise als separates CTAN-Paket
%  installiert). Ohne es kennt babel die Option "ngerman"
%  nicht und bricht den Build ab.
%
%  Alternativen und warum sie hier nicht greifen:
%   - \usepackage{ngermanb}   --> erwartet ngermanb.ldf, fehlt
%   - \usepackage[german]{babel} --> braucht german.ldf, fehlt
%   - polyglossia  --> nur fuer XeLaTeX/LuaLaTeX, wir nutzen pdfLaTeX
%
%  Was wir stattdessen tun:
%   - inputenc + fontenc:  korrekte UTF-8-Eingabe und T1-Ausgabe
%     (Umlaute in Fliesstextbloecken funktionieren damit einwandfrei)
%   - csquotes:  deutsche Anführungszeichen per \enquote{}
%   - Listings: extendedchars=false (verhindert UTF-8-Fehler
%     INNERHALB von lstlisting-Umgebungen; Kommentare dort
%     deshalb in ASCII-Umschreibung)
%   - Silbentrennung laeuft ohne babel auf Englisch; das ist
%     bei einem technischen Dokument mit vielen Code-Bezeichnern
%     kaum sichtbar.
%
%  Sobald babel-german auf dem System verfuegbar ist, reicht
%  es, die beiden auskommentierten Zeilen unten zu aktivieren
%  und die "\usepackage{csquotes}"-Zeile zu behalten.
%
%  Zum Aktivieren (wenn babel-german installiert):
%    \usepackage[ngerman]{babel}
%    \usepackage{csquotes}   % arbeitet dann mit babel zusammen
% ============================================================

% -- Zeichenkodierung ----------------------------------------
% inputenc: Eingabedatei wird als UTF-8 gelesen.
%   Ohne dieses Paket muss jeder Umlaut als Escape-Sequenz
%   geschrieben werden (\"a statt ae, etc.).
\usepackage[utf8]{inputenc}

% fontenc: Ausgabe-Zeichensatz T1 (8-bit).
%   T1 enthaelt alle westeuropaeischen Zeichen als echte
%   Glyphen -- Silbentrennung an Umlauten funktioniert damit
%   korrekt. Ohne T1 wuerde z.B. "Oe" als zwei Zeichen
%   behandelt und nicht sauber getrennt.
\usepackage[T1]{fontenc}

% csquotes: typografisch korrekte Anführungszeichen.
%   \enquote{Text} erzeugt "Text" (je nach Locale).
%   Ohne babel-Kopplung nutzen wir den Stil manuell:
%   \MakeOuterQuote{"} wäre eine Alternative, aber
%   \enquote ist sicherer und semantisch klarer.
\usepackage{csquotes}

% -- Schriften -----------------------------------------------
% palatino: Serifenschrift (Fliesstext).
%   Besser lesbar als Computer Modern auf Bildschirm und
%   Ausdruck; wird von Qt-Dokumentationen aehnlich verwendet.
\usepackage{palatino}

% courier: Monospace-Schrift fuer Code-Listings.
%   Ist auf praktisch jedem TeX-System vorhanden und
%   unterstuetzt den T1-Zeichensatz vollstaendig.
\usepackage{courier}

% helvet: serifenlose Schrift (Helvetica) fuer Ueberschriften,
%   skaliert auf 92 % der normalen Groesse (wirkt proportionaler).
\usepackage[scaled=0.92]{helvet}

% -- Seitenlayout --------------------------------------------
% geometry: Seitenraender fein einstellen.
%   inner/outer statt left/right fuer zweiseitigen Druck:
%   "inner" ist der Bundsteg (nahe am Heftrand), "outer" der
%   Aussensteg. headheight=15pt verhindert eine fancyhdr-Warnung.
\usepackage[
  top=2.8cm, bottom=3.2cm,
  inner=3.0cm, outer=2.4cm,
  headheight=15pt
]{geometry}

% fancyhdr: individuelle Kopf- und Fusszeilen.
%   Standardmaessig zeigt LaTeX nur Seitennummern;
%   fancyhdr erlaubt Kapiteltitel + Icon in der Kopfzeile.
\usepackage{fancyhdr}

% -- Farben --------------------------------------------------
% xcolor mit dvipsnames: vordefinierte Farbnamen (z.B. BrickRed)
%   und table-Option fuer gefaerbte Tabellenzeilen.
\usepackage[dvipsnames,table]{xcolor}

% -- Schriften (Ergaenzung) ----------------------------------
% microtype: verbessert den Textumbruch durch optischen
%   Randausgleich und Zeichenabstandsanpassung (nur pdfLaTeX).
%   Rein typografisch, kein sichtbarer Effekt auf Struktur.
\usepackage{microtype}

% -- Grafik --------------------------------------------------
% graphicx: \includegraphics fuer PNG/PDF/JPG.
%   Benoetigt fuer das Icon in Titelseite und Kopfzeile.
\usepackage{graphicx}

% -- Code-Listings -------------------------------------------
% listings: typsetzt Quellcode mit Syntax-Highlighting.
%   Eigene Stile werden in macros.tex definiert.
%   Hinweis: extendedchars=false ist in den Stilen gesetzt,
%   weil listings intern keine echte UTF-8-Unterstuetzung hat
%   und bei Umlauten in Code-Bloecken sonst abbricht.
\usepackage{listings}

% -- Boxen (Hinweise, Warnungen, Qt-Hintergrund) -------------
% tcolorbox mit skins und breakable:
%   "skins" erlaubt erweiterte Rahmengestaltung,
%   "breakable" erlaubt Boxen ueber Seitenumbrueche hinweg.
\usepackage{tcolorbox}
\tcbuselibrary{skins,breakable}

% -- Tabellen ------------------------------------------------
% booktabs: professionelle horizontale Linien (\toprule etc.)
%   ohne vertikale Linien -- entspricht typografischen Regeln.
\usepackage{booktabs}

% longtable: Tabellen ueber mehrere Seiten.
%   Benoetigt fuer die Kuerzel- und Klassen-Referenz im Anhang.
\usepackage{longtable}

% array: erweiterte Spalten-Definitionen in Tabellen
%   (z.B. feste Spaltenbreite mit p{3cm}).
\usepackage{array}

% -- Listen --------------------------------------------------
% enumitem: feinere Kontrolle ueber Listenabstaende.
%   noitemsep reduziert den Abstand zwischen Listenpunkten
%   fuer kompaktere Darstellung.
\usepackage{enumitem}

% -- Absatzgestaltung ----------------------------------------
% parskip: ersetzt Einrueckung durch halben Absatzabstand.
%   Passt besser zu einem technischen Dokument als klassische
%   Einrueckung; kombiniert mit scrreprt's parskip=half.
\usepackage{parskip}

% -- Index ---------------------------------------------------
% makeidx: erzeugt einen Stichwortindex am Dokumentende.
%   \index{Begriff} im Text, \printindex am Ende.
\usepackage{makeidx}
\makeindex

% -- Hyperlinks und PDF-Metadaten ----------------------------
% hyperref: macht Querverweise, TOC-Eintraege und URLs
\usepackage[
  colorlinks = true,
  linkcolor  = tvpurple2,%linkblue,
  urlcolor   = tvpurple2,%linkblue,
  citecolor  = gray,
  pdftitle   = {TilVista Entwicklerdokumentation},
  pdfauthor  = {Ludwig Armin}
]{hyperref}

% HINWEIS: Diese Datei wird per \input{} eingebunden,
% nicht als eigenstaendiges Dokument kompiliert.
